\section{Background and Related Work}
\label{sec:background}

This section introduces the four strands of work that the paper connects.
Section~\ref{sec:bg-cww} summarises Computing with Words and the \Perc{}.
Section~\ref{sec:bg-judges} reviews how large language models are used to judge mental-health replies.
Section~\ref{sec:bg-disagreement} presents the view that rater disagreement is information rather than noise.
Section~\ref{sec:bg-cities} places the work in research on cognitive cities and in the practice of public crisis services.

\subsection{Computing with Words and the Perceptual Computer}
\label{sec:bg-cww}

\emph{Computing with Words} (CWW) is a methodology in which the objects of computation are words and propositions drawn
from natural language rather than numbers \citep{zadeh1996fuzzy}.
Zadeh motivates CWW by two observations: much of the information people use is too imprecise to justify numbers, and
tolerance for imprecision can be exploited to obtain tractable and robust solutions.
In this view, a word is modelled as a fuzzy set: a set whose elements belong to it to a degree between zero and one.

The \emph{\Perc{}} (Per-C) is an architecture that makes CWW operational for subjective judgements
\citep{mendel2001perceptual,mendel2010perceptual}.
It consists of three components.
An \emph{encoder} transforms words into fuzzy-set models that are collected in a \emph{codebook}.
A \emph{CWW engine} computes with these models, for example by aggregation or by if--then rules.
A \emph{decoder} maps the engine output back to a word, a rank or a class that a person can act on.
The design goal is that the output of the Per-C should be useful to, and interpretable by, a human decision maker.

Mendel argues that word models should be \emph{interval type-2 fuzzy sets} (IT2 FSs) because words mean different
things to different people \citep{mendel2007computing}.
A type-1 fuzzy set assigns each element a single membership grade.
An IT2 FS instead assigns each element an interval of grades, bounded below by a \emph{lower membership function}
(LMF) and above by an \emph{upper membership function} (UMF).
The region between the two is the \emph{footprint of uncertainty} (FOU); its width expresses how much people disagree
about the word.
IT2 FSs are therefore a natural representation when the uncertainty of interest is disagreement between raters.

Mendel also proposes three principles for validating CWW \citep{mendel2007computing}.
Following Turing, a CWW system should be judged by whether its output is indistinguishable from that of human experts.
Following Popper, a word model should be open to refutation by data about how people use the word, which favours models
that are estimated from such data and that represent its uncertainty.
Following Occam, the simplest word model and engine that pass the validation should be preferred.
A later forum of leading CWW researchers showed that the field holds several complementary views of what CWW means
\citep{mendel2010means}.
This paper adopts Mendel's engineering view: CWW as a Per-C that supports a human judgement.

The CWW engines used in this paper come from the Per-C literature.
The \emph{ordered weighted average} (OWA) aggregates values after sorting them, so that its weights express an
attitude between pessimism (emphasising the smallest values) and optimism \citep{yager1988owa}.
The \emph{fuzzy weighted average} (FWA) generalises the weighted average to fuzzy inputs and weights by computing it on
each $\alpha$-cut with the Karnik--Mendel algorithms \citep{liu2008fwa}.
The \emph{linguistic weighted average} (LWA) extends the FWA to IT2 inputs and IT2 weights \citep{wu2007lwa}.
Decoders typically compare the engine output with each codebook word by the \emph{Jaccard similarity} of IT2 FSs, and
summarise an IT2 FS by its \emph{centroid}, an interval computed with the (Enhanced) Karnik--Mendel algorithms
\citep{karnik2001centroid,wu2009ekm,wu2009comparative}.
A complementary decoder is the \emph{2-tuple} representation, which reports a word together with a signed offset that
keeps the information lost by rounding to the nearest word \citep{herrera2000twotuple}.

\subsection{LLMs as Judges of Mental-Health Replies}
\label{sec:bg-judges}

\emph{LLM-as-a-judge} denotes the use of a large language model (LLM) to rate the output of another system, usually
in place of human raters \citep{zheng2023judging}.
Because a single judge can be biased and inconsistent, several studies aggregate a \emph{jury} of different LLM
judges, typically by averaging or voting \citep{verga2024juries}.
In mental health, the judge's verdict effectively decides which replies count as safe, so its reliability is a safety
question rather than a matter of benchmark convenience.

CounselBench provides large-scale expert evaluations of answers by LLMs and by online human therapists to real patient
questions, rated by mental-health professionals on six clinically grounded dimensions \citep{li2026counselbench}.
Its authors report that LLM judges systematically overrate model answers and overlook safety concerns that human
experts identify, most notably unauthorised medical advice.
Badawi et al.\ study, with large-scale benchmarks, when LLM judges can be trusted to evaluate mental-health replies,
including the agreement between LLM and human ratings \citep{badawi2026trust}.

\emph{Between Help and Harm} (BHH) addresses crisis handling directly \citep{arnaizrodriguez2026help}.
It introduces a clinically informed taxonomy of six crisis categories, a validation set of user messages labelled by
human annotators and by LLMs, and an expert-designed protocol that rates the appropriateness of a reply on a five-point
scale from harmful to fully appropriate.
Response quality in BHH is measured by an LLM-as-a-judge that the authors validated against human expert feedback.
The present paper uses the published BHH validation data and the stored judge outputs, so that no new model calls are
needed.

Three further studies define the context of this paper without being analysed in it.
MentalHealthBench is an open benchmark of realistic mental-health conversations built with licensed mental-health
experts and scored against expert-written rubrics \citep{openai2026mentalhealthbench}.
SIM-VAIL audits chatbots in multi-turn conversations with simulated users and shows that concerning behaviour
accumulates over turns \citep{weilnhammer2026simvail}.
The Verily Mental Health Guardrail is a dedicated crisis detector that prioritises sensitivity in order to minimise
missed psychiatric crises \citep{nelson2026guardrail}.
Together these studies show that safety evaluation is moving towards expert-grounded, multi-rater and multi-turn
designs.

A common pattern runs through this literature.
Whenever several judges, several runs or several experts produce ratings, the ratings are reduced to a single number
or label, usually a mean or a majority vote, before they are compared with a reference.
The spread among the raters, which may indicate that a reply is ambiguous or risky, is discarded at that step.
This paper asks what is gained when the spread is kept.

\subsection{Disagreement as Signal}
\label{sec:bg-disagreement}

Research on annotation increasingly treats disagreement between raters as information rather than as noise.
Uma et al.\ survey methods that learn from disagreement instead of resolving it into a single aggregated gold label
\citep{uma2021learning}.
Plank argues that human label variation is often genuine and that evaluation should account for it rather than assume a
single ground truth \citep{plank2022problem}.
Crisis assessment is a domain where such variation is expected: whether a message signals a crisis can depend on
context that a rater does not have.

The \emph{Interval Agreement Approach} (IAA) offers a principled way to keep disagreement inside a fuzzy set
\citep{wagner2015iaa}.
Each rater contributes an interval, and the membership of a value is the proportion of intervals that contain it.
The IAA preserves both the variation between raters and the uncertainty within one rater, without assuming a
parametric shape.
Section~\ref{sec:concept} uses the IAA to encode panels of LLM judges, where the runs of one judge play the role of
within-rater uncertainty.

\subsection{Cognitive Cities and Public Crisis Services}
\label{sec:bg-cities}

\emph{Cognitive cities} extend the smart-city idea from efficient infrastructure to urban systems that learn and
interact with their citizens \citep{finger2016cognitive}.
In this line of work, soft computing and CWW are proposed as a way to let public services reason with the imprecise,
linguistic information that citizens provide.
D'Onofrio, M\"uller and Portmann describe a fuzzy reasoning process for conversational agents in cognitive cities,
combining fuzzy inference with fuzzy cognitive maps \citep{donofrio2019fuzzy}.
The present paper follows this direction for one specific public service: a crisis line that uses a chatbot.

Public crisis services have begun to deploy conversational AI.
In October 2026, the Swiss helpline Die Dargebotene Hand (telephone number 143) began to offer an AI chatbot in its
online chat for people who cannot reach a human counsellor \citep{srf2026dargebotene}.
The Dutch suicide-prevention line 113 Zelfmoordpreventie follows a different design: it uses AI only in chat
conversations to support its counsellors, and it states that a helpline conversation always takes place with a human
\citep{zelfmoordpreventie113hulplijn,vu2024aitool}.
Both designs need a decision about when an automated reply is not enough and a person in crisis should be handed over
to a human.
The escalation decoder proposed in Section~\ref{sec:concept} is a concept for making that decision transparent and
sensitive to disagreement.
